% Vietnamese translation for OI Public Library
% Source-File: IOI中国国家候选队论文/2005/栗师/栗师.ppt
\documentclass[11pt]{article}
\usepackage{oipl-vn}

\title{Bản trình chiếu: Ghi chú về tổng khoảng cách ngắn nhất}
\author{Li Shi}
\date{2005}

\begin{document}
\maketitle

\origtitle{Slide companion for Li Shi's shortest-distance presentation}
\sourcepath{IOI China candidate team papers / 2005 / Li Shi / Li Shi.ppt}

\section{Phạm vi bản dịch}

Tệp gốc chỉ có PowerPoint. Phần chữ đọc được lặp lại công thức
\[
  d[i]=\min\{\operatorname{dis}[i,j]\},
  \qquad d[1]+d[2]+\cdots+d[n],
\]
cùng một số đối tượng hình ảnh. Nội dung dưới đây ghi lại mô hình thuật toán
có thể xác định từ các công thức đó.

\section{Mô hình}

Với mỗi điểm hoặc đỉnh \(i\), đại lượng \(d[i]\) là khoảng cách tốt nhất từ
\(i\) tới một tập đích hoặc một điểm được chọn. Mục tiêu là tối ưu tổng
\(\sum_i d[i]\). Đây là dạng bài thường gặp trong chọn trung tâm, chọn trạm,
hoặc gom nhóm trên đồ thị.

\section{Hướng giải}

Nếu \(\operatorname{dis}[i,j]\) là khoảng cách trên đồ thị, bước đầu là tính
khoảng cách ngắn nhất bằng BFS, Dijkstra, Floyd hoặc biến thể theo cấu trúc
đề. Sau đó bài toán trở thành chọn \(j\) hoặc một tập các \(j\) sao cho tổng
khoảng cách nhỏ nhất. Khi số đỉnh lớn, cần khai thác tính chất đặc biệt thay
vì thử mọi tập.

\section{Ghi chú}

Do phần văn bản chính của slide không đọc được đầy đủ, bản dịch này giữ vai
trò companion note và không mở rộng thành một lời giải cụ thể hơn dữ liệu
khôi phục được.

\end{document}
